\lyseis{1.2}{Πρόσθεση, αφαίρεση, πολλαπλασιασμός}
\sol{1} Στο $24+15=39$: προσθετέοι και άθροισμα. Στο $50-18=32$: μειωτέος, αφαιρετέος και διαφορά. Στο $6\cdot 9=54$: παράγοντες και γινόμενο.
\sol{2} α) αντιμεταθετική της πρόσθεσης \quad β) προσεταιριστική της πρόσθεσης \quad γ) επιμεριστική του πολλαπλασιασμού ως προς την πρόσθεση \quad δ) το 1 δεν αλλάζει το γινόμενο \quad ε) ο πολλαπλασιασμός με το 0 δίνει 0 \quad στ) το 0 δεν αλλάζει το άθροισμα \quad ζ) αντιμεταθετική του πολλαπλασιασμού \quad η) επιμεριστική ως προς την αφαίρεση.
\sol{3} α) 630 \quad β) 20.500 \quad γ) 38.000 \quad δ) 700.000.
\sol{4} α) $71-26=45$ \quad β) $100-38=62$ \quad γ) $45+19=64$ \quad δ) $90-57=33$. Έλεγχος, π.χ. στο δ: $90-33=57$.
\sol{5} α) $2\cdot 50\cdot 37=100\cdot 37=3.700$ \quad β) $25\cdot 4\cdot 13=100\cdot 13=1.300$ \quad γ) $(16+84)+(39+61)=100+100=200$ \quad δ) $5\cdot 20\cdot 49=100\cdot 49=4.900$.
\sol{6} α) $56\cdot(8+2)=56\cdot 10=560$ \quad β) $(37+63)\cdot 64=100\cdot 64=6.400$ \quad γ) $91\cdot(15-5)=91\cdot 10=910$ \quad δ) $47\cdot(100-1)=4.700-47=4.653$ \quad ε) $24\cdot(100+1)=2.400+24=2.424$.
\sol{7} α) Λάθος: το $5-12$ δεν υπάρχει στους φυσικούς. β) Λάθος: $(30-10)-5=15$, ενώ $30-(10-5)=25$. γ) Λάθος: $4\cdot(6+1)=28$, ενώ $4\cdot 6+1=25$. δ) Λάθος: $0\cdot 9=0$. ε) Σωστό. στ) Σωστό.
\sol{8} α) 300 \quad β) $(117+83)+46=246$ \quad γ) 528 \quad δ) 3.486 (έλεγχος: $2.518+3.486=6.004$) \quad ε) $90-25=65$.
\sol{9} α) $(8+5)+2=15$ και $8+(5+2)=15$. β) $3\cdot(4\cdot 5)=60$ και $(3\cdot 4)\cdot 5=60$. Είναι η προσεταιριστική ιδιότητα, της πρόσθεσης στο α) και του πολλαπλασιασμού στο β).
\sol{10} $98\cdot 5=(100-2)\cdot 5=500-10=490$ τοστ.
\sol{11} Αφετηρία ως γέφυρα: $47$ χλμ. Γέφυρα ως βενζινάδικο: $112-47=65$ χλμ. Βενζινάδικο ως χωριό: $175-112=63$ χλμ. Χωριό ως θάλασσα: $230-175=55$ χλμ. Βενζινάδικο ως θάλασσα: $230-112=118$ χλμ. Έλεγχος: $63+55=118$ και $47+65+63+55=230$.
\sol{12} Μειωτέος είναι τα αυτοκόλλητα που είχε: $Μ=17+38=55$. Έλεγχος: $55-17=38$.
\sol{13} Στην αφαίρεση $82-47$ το 2 στις μονάδες είναι μικρότερο από το 7. Δεν επιτρέπεται να γυρίσουμε τους αριθμούς ανάποδα: δανειζόμαστε μία δεκάδα και κάνουμε $12-7=5$, και μετά $7-4=3$. Σωστό: 35. Έλεγχος με $Μ=Α+\Delta$: $47+35=82$. Με το 45 του Ηλία βγαίνει $47+45=92$, που δεν είναι 82.
\sol{14} $48\cdot 25=12\cdot 4\cdot 25=12\cdot 100=1.200$ πλακάκια.
\sol{15} $13\cdot 8=(10+3)\cdot 8=10\cdot 8+3\cdot 8=80+24=104$ τ.μ. Ένα ορθογώνιο $10\times 8$ και ένα $3\times 8$, κολλημένα.
\sol{16} $3\cdot 99=3\cdot(100-1)=300-3=297$ €.
\sol{17} Τα ζευγάρια είναι 30 και το καθένα κάνει 61 ($1+60$, $2+59$, \dots, $30+31$), άρα $30\cdot 61=1.830$ κούτσουρα. Για το $1+\dots+8$: 4 ζευγάρια και καθένα κάνει 9, άρα $4\cdot 9=36$. Έλεγχος με πρόσθεση: $1+2+3+4+5+6+7+8=36$.
\sol{18} Οι αριθμοί είναι 13 και 7: $13+7=20$ και $13-7=6$.
\sol{19} Πέντε κουτιά με 0 κουλουράκια έχουν συνολικά $0+0+0+0+0=0$. Γενικά $5\cdot 0=0+0+0+0+0=0$.
\sol{20} Το ορθογώνιο με 3 σειρές των 5 έχει 15 τετράγωνα. Αν το γυρίσεις στο πλάι, βλέπεις 5 σειρές των 3, και τα τετράγωνα είναι πάλι 15. Η εικόνα είναι διαφορετική (3 ομάδες των 5 ή 5 ομάδες των 3), το αποτέλεσμα ίδιο.
\sol{21} $(20-8)-4=12-4=8$ και $20-(8-4)=20-4=16$. Οι τιμές διαφέρουν, άρα η αφαίρεση δεν είναι προσεταιριστική: έχει σημασία ποιες δύο αφαιρέσεις γίνονται πρώτες.
\sol{22} $3\cdot 3=9$, $9\cdot 3=27$, $27\cdot 3=81$, $81\cdot 3=243$, $243\cdot 3=729$. Έχει 3 ψηφία. Πιο σύντομα: $3^6$, «το τρία στην έκτη».
